\documentclass[11pt]{amsart}

\usepackage[T1]{fontenc}
\usepackage{lmodern}
\usepackage{microtype}
\usepackage{mathtools,amssymb,amsthm}
\usepackage[hidelinks]{hyperref}

\hypersetup{
  pdftitle={Constrained Ramsey Numbers for the Loose 3-Uniform Path}
}

\newtheorem{theorem}{Theorem}
\newtheorem{lemma}[theorem]{Lemma}

\newcommand{\Lpath}{\mathcal L}
\newcommand{\Z}{\mathbb Z}

\title[Constrained Ramsey numbers for the loose path]{Constrained Ramsey Numbers for the Loose 3-Uniform Path}
\date{}
\subjclass[2020]{05C55, 05D10}
\keywords{constrained Ramsey number, rainbow Ramsey number,
3-uniform hypergraph, loose path}

\begin{document}

\begin{abstract}
Let $\Lpath$ be the loose $3$-uniform path of length three. We prove that
\[
  f(H,\Lpath)=\max\{R_2(H),7\}
\]
for every $3$-graph $H$ with at least two edges. The new case is the order
boundary $R_2(H)=|V(H)|\ge 7$, and therefore the result answers
Problem~7.2 of Li and Wang. The proof uses their structural theorem for
rainbow-$\Lpath$-free colorings, their transfer argument for the second
structural case, and a necessary structural condition on $3$-graphs
attaining the order bound for their two-color Ramsey number.
\end{abstract}

\maketitle

\section{Introduction}

Let
\[
\Lpath=
\bigl\{
  \{v_1,v_2,v_3\},
  \{v_3,v_4,v_5\},
  \{v_5,v_6,v_7\}
\bigr\}
\]
be the loose $3$-uniform path of length three. For a $3$-graph $H$, let
$R_2(H)$ be the least $n$ such that every two-coloring of $K_n^{(3)}$
contains a monochromatic copy of $H$. For $3$-graphs $H$ and $G$, let
$f(H,G)$ be the least $n$ such that every edge-coloring of $K_n^{(3)}$,
with an arbitrary number of colors, contains either a monochromatic $H$
or a rainbow $G$. Every copy of $H$ in $K_n^{(3)}$ occupies $|V(H)|$
vertices, so $K_n^{(3)}$ contains no copy of $H$ at all when
$n<|V(H)|$; hence we have the \emph{order bound}
\[
  R_2(H)\ge|V(H)|.
\]
This inequality is used twice in the proof of Theorem~\ref{thm:main}.

Liu proved $f(H,\Lpath)=R_2(H)$ when
$R_2(H)\ge\max\{|V(H)|+3,10\}$~\cite{Liu}; we have not been able to
consult \cite{Liu} itself and quote this statement from the account of it
in \cite[Theorem~1.7]{LiWang}. Li and Wang improved this to
\[
  R_2(H)\ge\max\{|V(H)|+1,7\}
\]
and asked whether equality still holds at the remaining boundary
$R_2(H)=|V(H)|\ge7$~\cite[Theorem~1.10 and Problem~7.2]{LiWang}.
We answer their problem affirmatively and obtain the complete formula.

\begin{theorem}\label{thm:main}
For every $3$-graph $H$ with at least two edges,
\[
  f(H,\Lpath)=\max\{R_2(H),7\}.
\]
\end{theorem}

For a $3$-graph with one edge, one trivially has
$f(H,\Lpath)=R_2(H)=|V(H)|$. Thus Theorem~\ref{thm:main} determines
$f(H,\Lpath)$ for every nonempty $3$-graph $H$.

We use the following theorem of Li and Wang~\cite[Theorem~1.5]{LiWang}.

\begin{theorem}[Li--Wang]\label{thm:LW}
Let $n\ge7$, and let $c$ be an edge-coloring of $K_n^{(3)}$ using at
least three colors and containing no rainbow $\Lpath$. Then at least one
of the following holds:
\begin{enumerate}
\item there is a vertex $u$ such that $K_n^{(3)}-u$ is monochromatic;
\item there are an edge $e$ and a color $\alpha\ne c(e)$ such that every
edge $f\ne e$ with $c(f)\ne\alpha$ satisfies $|e\cap f|=2$.
\end{enumerate}
\end{theorem}

\section{The order boundary}

A $3$-graph $H$ is a \emph{cone} over a graph $J$ if
\[
  V(H)=\{x\}\cup V(J),
  \qquad
  E(H)=\bigl\{\{x,p,q\}:pq\in E(J)\bigr\}.
\]
The vertex $x$ is the cone vertex and $J$ is its link graph. We write
$K_a+K_b$ for the disjoint union of $K_a$ and $K_b$.

\begin{lemma}[Boundary structure]\label{lem:boundary}
Let $H$ be a $3$-graph on $n\ge7$ vertices with $R_2(H)=n$. Then either
$H$ has an isolated vertex or $H$ is a cone over a graph $J$ with
$\delta(J)=1$.
\end{lemma}

\begin{proof}
Fix $z\in V(K_n^{(3)})$, and color a triple red if it contains $z$ and
blue otherwise. Since every copy of $H$ in $K_n^{(3)}$ is spanning, a
red copy makes $H$ a cone and a blue copy gives an isolated vertex. Thus,
if $H$ has no isolated vertex, it is a cone over a graph $J$ on
$m=n-1$ vertices, and $\delta(J)\ge1$.

Suppose that $\delta(J)\ge2$. For $A\subseteq V(K_n^{(3)})$, color each
triple $T$ by $|T\cap A|\pmod 2$. If $|A|=a$, then the two links at a
vertex are
\[
\begin{cases}
K_{a-1,n-a}\ \text{and}\ K_{a-1}+K_{n-a},&y\in A,\\
K_{a,n-a-1}\ \text{and}\ K_a+K_{n-a-1},&y\notin A.
\end{cases}
\]
A monochromatic copy of the cone would place a spanning copy of $J$ in
one of these links.

First assume $n\ge8$, so $m\ge7$. With $a=2$, either the coloring avoids
$H$, or $J$ is a spanning subgraph of one of
\[
  K_{1,m-1},\quad K_1+K_{m-1},\quad
  K_{2,m-2},\quad K_2+K_{m-2}.
\]
The minimum-degree condition excludes all but $K_{2,m-2}$ and then forces
$J\cong K_{2,m-2}$. Taking $a=4$, the possible links are
\[
  K_{3,m-3},\quad K_3+K_{m-3},\quad
  K_{4,m-4},\quad K_4+K_{m-4}.
\]
None contains a spanning $K_{2,m-2}$: the disconnected links cannot
contain a connected spanning graph, while a connected bipartite graph has
a unique bipartition and the complete bipartite links have different part
sizes. Hence one of the two parity colorings avoids $H$, contradicting
$R_2(H)=n$.

It remains to treat $n=7$. The $a=2$ coloring either avoids $H$ or forces
$J\cong K_{2,4}$. In the latter case, work on $\Z_7$ and set
\[
  D_1=\{0,1,3\},\qquad D_2=\{0,2,3\},\qquad
  \mathcal B_i=\{t+D_i:t\in\Z_7\}.
\]
The ordered nonzero differences of each $D_i$ are all six nonzero
residues, so each $\mathcal B_i$ is a Steiner triple system. The two
systems are disjoint: the translates of $D_1$ containing $0$ are
$\{0,1,3\}$, $\{0,2,6\}$, and $\{0,4,5\}$, none of which equals $D_2$.
Color the fourteen blocks in $\mathcal B_1\cup\mathcal B_2$ red and all
other triples blue. At every vertex, the red link is the union of two
edge-disjoint perfect matchings, hence a $6$-cycle; the blue link is its
$3$-regular complement in $K_6$. Neither link contains a spanning
$K_{2,4}$: the red link has six edges and the blue link has maximum degree
three. This again contradicts $R_2(H)=7$.

Therefore $\delta(J)=1$.
\end{proof}

The next lemma is not new: it is the argument Li and Wang give for
case~(ii), which is part~(2) above, in their proof of
\cite[Theorem~1.10]{LiWang}, where the
recoloring is set up through the two isomorphic spanning subgraphs
$F_1'=\{v_iv_jv_\ell:i,j\in[3],\ \ell\notin[3]\}\cup\{v_1v_2v_3\}$ and
$F_2'$ defined in the same way from $\{v_4,v_5,v_6\}$; the permutation
$\pi$ below is their isomorphism $F_1'\to F_2'$. Their proof is stated
for $n=R_2(H)$, and it extends without change to $n\ge R_2(H)$, which is
the only form we need. We include it for completeness and because the
proof of Theorem~\ref{thm:main} applies it at $N=\max\{R_2(H),7\}$.

\begin{lemma}[Transfer; Li and Wang \cite{LiWang}]\label{lem:transfer}
Let $n\ge7$, let $R_2(H)\le n$, and suppose that a rainbow-$\Lpath$-free
coloring of $K_n^{(3)}$ satisfies part~(2) of
Theorem~\ref{thm:LW}. Then it contains a monochromatic copy of $H$.
\end{lemma}

\begin{proof}
Relabel so that $e=\{1,2,3\}$, choose distinct vertices $4,5,6$ outside
$e$, and let
\[
  \pi=(1\ 4)(2\ 5)(3\ 6)
\]
fix all remaining vertices. Let $S$ be the set of edges whose color is
not $\alpha$. If $f=e$, then $\pi(f)$ is disjoint from $e$; if $f\ne e$,
then $|f\cap e|=2$ and $|\pi(f)\cap e|\le1$. In either case
$\pi(f)\ne e$, so the contrapositive of part~(2) gives
$c(\pi(f))=\alpha$.

Recolor the edges in $S$ blue and all other edges red. Since
$R_2(H)\le n$, this two-coloring contains a monochromatic $H$. A red copy
is already $\alpha$-monochromatic in the original coloring, while the
image under $\pi$ of a blue copy is $\alpha$-monochromatic.
\end{proof}

\section{Proof of the formula}

\begin{proof}[Proof of Theorem~\ref{thm:main}]
Write $r=R_2(H)$ and $N=\max\{r,7\}$. A two-coloring of
$K_{r-1}^{(3)}$ without a monochromatic $H$ contains no rainbow
$\Lpath$, so $f(H,\Lpath)\ge r$. Also, coloring all triples of
$K_6^{(3)}$ differently gives neither a monochromatic $H$ nor a rainbow
$\Lpath$. Hence $f(H,\Lpath)\ge N$.

For the reverse inequality, let $c$ be a rainbow-$\Lpath$-free coloring
of $K_N^{(3)}$. If at most two colors occur, then $c$ contains a
monochromatic $H$ because $N\ge r$. Assume that at least three colors
occur and apply Theorem~\ref{thm:LW}.

If part~(2) holds, Lemma~\ref{lem:transfer} gives a monochromatic $H$.
Suppose part~(1) holds, and let $u$ be such that all triples in
$W=V(K_N^{(3)})\setminus\{u\}$ have one color, say $\alpha$.

If $r\le6$, then $N=7$ and, by the order bound, $|V(H)|\le r\le6$, so
the monochromatic $K_6^{(3)}$ on $W$ contains $H$. Now let $r\ge7$, so
$N=r$. If $r\ge|V(H)|+1$, the same conclusion is immediate; this is the
step Li and Wang use for case~(i), that is part~(1), in their proof of
\cite[Theorem~1.10]{LiWang}. By the order bound it remains only to
consider
\[
  r=|V(H)|\ge7.
\]
By Lemma~\ref{lem:boundary}, either $H$ has an isolated vertex or $H$ is
a cone over a graph $J$ with $\delta(J)=1$. In the first case, map an
isolated vertex to $u$ and all remaining vertices bijectively to $W$.
This gives an $\alpha$-monochromatic copy of $H$.

In the second case, let $x$ be the cone vertex, choose a leaf $p$ of $J$,
and let $q$ be its unique neighbor. If
$c(\{u,y,w\})=\alpha$ for some distinct $y,w\in W$, map
$x,p,q$ to $y,u,w$, respectively, and extend bijectively. The edge
corresponding to $pq$ maps to $\{u,y,w\}$, and every other edge maps into
$W$; hence the resulting copy of $H$ is $\alpha$-monochromatic.

We may therefore assume that every triple through $u$ has color different
from $\alpha$. Color each pair $ab\in\binom{W}{2}$ by
$c(\{u,a,b\})$. If $ab$ and $cd$ are disjoint, choose distinct
$r',s'\in W\setminus\{a,b,c,d\}$. Then
\[
  \{u,a,b\},\qquad \{u,c,d\},\qquad \{c,r',s'\}
\]
form a copy of $\Lpath$, and the last edge has color $\alpha$. The first
two edges must therefore have the same color. The graph on
$\binom{W}{2}$ joining disjoint pairs is connected for $|W|\ge6$: two
distinct intersecting pairs have a common disjoint pair. Thus all triples
through $u$ have one common color, different from $\alpha$. Hence $c$
uses only two colors, a contradiction.

Every rainbow-$\Lpath$-free coloring of $K_N^{(3)}$ contains a
monochromatic $H$. Hence $f(H,\Lpath)\le N$, completing the proof.
\end{proof}

The case $R_2(H)=|V(H)|\ge7$ in the proof gives, in particular, an
affirmative answer to Problem~7.2 of Li and Wang.

\section*{Status and scope}

What is settled is Problem~7.2 of \cite{LiWang} as stated there: for every
$3$-graph $H$ with $R_2(H)=|V(H)|\ge7$ one has $f(H,\Lpath)=R_2(H)$.
Problem~7.2 imposes no restriction on the number of edges of $H$;
Theorem~\ref{thm:main} assumes $|E(H)|\ge2$, and the one-edge case is
disposed of in the paragraph following it, so the problem is answered in
full.

Two of the ingredients are due to Li and Wang.
Theorem~\ref{thm:LW} is \cite[Theorem~1.5]{LiWang}, and
Lemma~\ref{lem:transfer} is the case-(ii) argument in their proof of
\cite[Theorem~1.10]{LiWang}, in different notation and with $n=R_2(H)$
relaxed to $n\ge R_2(H)$; the step that disposes of $r\ge|V(H)|+1$ in
case~(1) is likewise theirs. We have not been able to consult the full
text of \cite{Liu}, so we do not know whether the recoloring-and-transfer
device of Lemma~\ref{lem:transfer} originates with \cite{LiWang} or
already with Liu; in either case it is not new here. What is new is
Lemma~\ref{lem:boundary} together with the treatment of case~(1) at
$r=|V(H)|$, which is the case Problem~7.2 asks about.

Lemma~\ref{lem:boundary} states a necessary condition and nothing more. It
does not classify the $3$-graphs with $R_2(H)=|V(H)|\ge7$, and no converse
is proved or claimed: a $3$-graph with an isolated vertex, or a cone over a
graph of minimum degree~$1$, need not attain the order bound.

The coloring used in the proof of Lemma~\ref{lem:boundary} for $n=7$ is a
pair of disjoint Steiner triple systems on $\Z_7$, that is, a pair of
disjoint Fano planes. Such a pair is classical and is not claimed here as a
new object; what is used is only that the two links it produces at a
vertex, a $6$-cycle and its $3$-regular complement in $K_6$, both fail to
contain a spanning $K_{2,4}$, which is what rules out a monochromatic cone
over $K_{2,4}$ in $K_7^{(3)}$.

\paragraph{Exact verification.}
The verifications for $n=7$ in the proof of Lemma~\ref{lem:boundary} are
finite and hand-sized, and each one is carried out from the data printed in
this paper: the six ordered nonzero differences of
each of $D_1=\{0,1,3\}$ and $D_2=\{0,2,3\}$, which give the Steiner property
of $\mathcal B_1$ and $\mathcal B_2$; the three translates of $D_1$ containing
$0$, which give their disjointness; and the two links at a vertex of the
resulting coloring of $K_7^{(3)}$, which give the failure of a spanning
$K_{2,4}$. The case $n\ge8$ of that lemma is not a finite computation: there
the two links produced by the parity coloring at a given value of $a$ are
described for general $m=n-1$, and the exclusions are made by the
minimum-degree, connectivity and bipartition arguments written out in the
proof rather than by inspecting a fixed list. No
program is needed and none accompanies this paper; a reader can redo every
step from the paper by hand.

We have found no answer to Problem~7.2 in
the literature: the only paper we found citing \cite{LiWang} is the
companion paper \cite{LiLiu}, on graphs and rainbow $P_5$, which does not
address $f(H,\Lpath)$; we found no journal version of \cite{LiWang}, and
its latest arXiv revision still records Problem~7.2 as open. Our search
is not exhaustive, and the result is offered subject to external
prior-art review.

\begin{thebibliography}{9}

\bibitem{Liu}
X.~Liu,
\emph{Constrained Ramsey numbers for the loose path, cycle and star},
Discrete Math.\ \textbf{347} (2024), no.~1, Paper No.~113697, 11 pp.
\href{https://doi.org/10.1016/j.disc.2023.113697}
{doi:10.1016/j.disc.2023.113697}.

\bibitem{LiWang}
X.~Li and R.~Wang,
\emph{Edge-colored $3$-uniform hypergraphs without rainbow paths of
length $3$ and its applications to Ramsey theory},
preprint, \href{https://arxiv.org/abs/2510.18246}{arXiv:2510.18246v2}
[math.CO] (2025);
\href{https://doi.org/10.48550/arXiv.2510.18246}
{doi:10.48550/arXiv.2510.18246}.

\bibitem{LiLiu}
X.~Li and X.~Liu,
\emph{Constrained Ramsey numbers for rainbow $P_5$},
preprint, \href{https://arxiv.org/abs/2510.18243}{arXiv:2510.18243v2}
[math.CO] (2025).

\end{thebibliography}

\end{document}